\documentclass{article}
\PassOptionsToPackage{numbers,sort&compress}{natbib}
\usepackage[sglblindworkshop, final]{neurips_2026}
\usepackage[utf8]{inputenc}
\usepackage[T1]{fontenc}
\usepackage{hyperref}
\hypersetup{hidelinks,hypertexnames=false}
\usepackage{amsmath,amssymb}
\usepackage{booktabs}
\usepackage{array}
\usepackage{graphicx}
\usepackage{microtype}
\usepackage{xcolor}
\usepackage{url}
\usepackage{float}
\newcommand{\pullPeerG}{0.06}
\newcommand{\sePeerG}{0.04}
\newcommand{\nPeerG}{96}
\newcommand{\pullToolG}{0.08}
\newcommand{\seToolG}{0.03}
\newcommand{\nToolG}{96}
\newcommand{\pullPeerN}{0.61}
\newcommand{\sePeerN}{0.09}
\newcommand{\nPeerN}{64}
\newcommand{\pullToolN}{0.61}
\newcommand{\seToolN}{0.08}
\newcommand{\nToolN}{64}
\newcommand{\premium}{-0.018}
\newcommand{\premiumse}{0.041}
\newcommand{\premiump}{0.69}
\newcommand{\premiumlo}{-0.10}
\newcommand{\premiumhi}{+0.06}
\newcommand{\premiumn}{192}
\newcommand{\premiumclusters}{16}
\newcommand{\ambig}{0.54}
\newcommand{\ambigse}{0.07}
\newcommand{\ambigp}{<0.001}
\newcommand{\triple}{0.01}
\newcommand{\triplese}{0.10}
\newcommand{\triplep}{0.91}
\newcommand{\triplelo}{-0.18}
\newcommand{\triplehi}{+0.20}
\newcommand{\facSev}{152}
\newcommand{\facPrior}{63}
\newcommand{\facRem}{0}
\newcommand{\facCoop}{-21}
\newcommand{\rangeInside}{54}
\newcommand{\rangeBelow}{45}
\newcommand{\rangeAbove}{1}
\newcommand{\DispG}{0.073}
\newcommand{\riPeerGp}{0.513}
\newcommand{\riPeerGest}{0.06}
\newcommand{\riToolGp}{0.022}
\newcommand{\riToolGest}{0.08}
\newcommand{\riPeerNp}{<0.001}
\newcommand{\riPeerNest}{0.61}
\newcommand{\riToolNp}{<0.001}
\newcommand{\riToolNest}{0.61}
\newcommand{\riClerNp}{0.690}
\newcommand{\riClerNest}{0.07}
\newcommand{\riPremGp}{0.667}
\newcommand{\riPremGest}{-0.02}
\newcommand{\riPremNp}{0.946}
\newcommand{\riPremNest}{-0.01}
\newcommand{\feAmbig}{0.54}
\newcommand{\feAmbigse}{0.07}
\newcommand{\feAgents}{6}
\newcommand{\driftNo}{0.01}
\newcommand{\driftNose}{0.06}
\newcommand{\reexNo}{-0.01}
\newcommand{\reexNose}{0.02}
\newcommand{\reexabsNo}{0.04}
\newcommand{\spreadNo}{0.029}
\newcommand{\driftOwn}{0.09}
\newcommand{\driftOwnse}{0.09}
\newcommand{\reexOwn}{0.01}
\newcommand{\reexOwnse}{0.02}
\newcommand{\reexabsOwn}{0.06}
\newcommand{\spreadOwn}{0.047}
\newcommand{\armgap}{0.006}
\newcommand{\armgapse}{0.011}
\newcommand{\mdlOpusPG}{0.06}
\newcommand{\mdlOpusPGn}{96}
\newcommand{\mdlOpusTG}{0.08}
\newcommand{\mdlOpusTGn}{96}
\newcommand{\mdlOpusPN}{0.61}
\newcommand{\mdlOpusPNn}{64}
\newcommand{\mdlOpusTN}{0.61}
\newcommand{\mdlOpusTNn}{64}
\newcommand{\twTool}{0.61}
\newcommand{\twPeer}{-0.01}
\newcommand{\twPeerse}{0.07}
\newcommand{\twPeerp}{0.93}
\newcommand{\twCler}{-0.54}
\newcommand{\twClerse}{0.11}
\newcommand{\twClerp}{<0.001}
\newcommand{\twN}{192}
\newcommand{\sdRatio}{10}
\newcommand{\NaNn}{62}
\newcommand{\NaNdev}{0.26}
\newcommand{\NaNsd}{0.286}
\newcommand{\NaNsev}{101}
\newcommand{\NaNprior}{74}
\newcommand{\NaGn}{48}
\newcommand{\NaGdev}{-0.20}
\newcommand{\NaGsd}{0.029}
\newcommand{\NaGsev}{162}
\newcommand{\NaGprior}{71}
\newcommand{\menPG}{52}
\newcommand{\menTG}{64}
\newcommand{\menPN}{6}
\newcommand{\menTN}{64}
\newcommand{\pullMention}{0.45}
\newcommand{\pullSilent}{0.67}
\newcommand{\exactAdopt}{22}
\newcommand{\pullClerN}{0.07}
\newcommand{\seClerN}{0.09}
\newcommand{\nClerN}{64}
\newcommand{\loClerN}{-0.11}
\newcommand{\hiClerN}{0.25}
\newcommand{\IntPN}{-0.05}
\newcommand{\IntTN}{-0.03}
\newcommand{\IntPG}{-0.12}
\newcommand{\menClerN}{31}
\newcommand{\RatGtop}{Proportionality}
\newcommand{\RatGtopv}{49}
\newcommand{\RatGsec}{Deterrence}
\newcommand{\RatGsecv}{29}
\newcommand{\RatNtop}{Deterrence}
\newcommand{\RatNtopv}{47}
\newcommand{\RatNsec}{Proportionality}
\newcommand{\RatNsecv}{30}
\newcommand{\yokedCells}{160}
\newcommand{\yokedMismatch}{0}
\newcommand{\worstImbalance}{0.010}
\newcommand{\gpOpus}{0.07}
\newcommand{\gxOpus}{0.54}
\newcommand{\gxseOpus}{0.07}
\newcommand{\prOpusG}{-0.02}
\newcommand{\prseOpusG}{0.04}
\newcommand{\prOpusN}{-0.01}
\newcommand{\prseOpusN}{0.07}
\newcommand{\gpSonnet}{0.24}
\newcommand{\gxSonnet}{0.63}
\newcommand{\gxseSonnet}{0.15}
\newcommand{\prSonnetG}{0.20}
\newcommand{\prseSonnetG}{0.06}
\newcommand{\prSonnetN}{0.07}
\newcommand{\prseSonnetN}{0.29}
\newcommand{\gpHaiku}{0.56}
\newcommand{\gxHaiku}{0.36}
\newcommand{\gxseHaiku}{0.06}
\newcommand{\prHaikuG}{0.21}
\newcommand{\prseHaikuG}{0.11}
\newcommand{\prHaikuN}{0.07}
\newcommand{\prseHaikuN}{0.12}
\newcommand{\prSmall}{0.14}
\newcommand{\prSmallse}{0.11}
\newcommand{\prSmallp}{0.26}
\newcommand{\prSmalln}{128}
\newcommand{\rbGCasese}{0.04}
\newcommand{\rbGCaselo}{-0.10}
\newcommand{\rbGCasehi}{+0.06}
\newcommand{\rbGAgentse}{0.04}
\newcommand{\rbGAgentlo}{-0.10}
\newcommand{\rbGAgenthi}{+0.07}
\newcommand{\rbGAssignse}{0.05}
\newcommand{\rbGAssignlo}{-0.12}
\newcommand{\rbGAssignhi}{+0.08}
\newcommand{\rbGType}{-0.018}
\newcommand{\rbGTypese}{0.042}
\newcommand{\rbNCasese}{0.07}
\newcommand{\rbNCaselo}{-0.14}
\newcommand{\rbNCasehi}{+0.13}
\newcommand{\rbNAgentse}{0.14}
\newcommand{\rbNAgentlo}{-0.28}
\newcommand{\rbNAgenthi}{+0.27}
\newcommand{\rbNAssignse}{0.11}
\newcommand{\rbNAssignlo}{-0.22}
\newcommand{\rbNAssignhi}{+0.21}
\newcommand{\rbNType}{-0.006}
\newcommand{\rbNTypese}{0.069}
\newcommand{\plSd}{0.06}
\newcommand{\plPctile}{0.10}
\newcommand{\plMax}{0.11}
\newcommand{\plN}{40}
\newcommand{\dispRatio}{10}
\newcommand{\dispUnThree}{0.29}
\newcommand{\pullPeerBare}{0.88}
\newcommand{\sePeerBare}{0.06}
\newcommand{\nPeerBare}{64}
\newcommand{\pullToolBare}{0.55}
\newcommand{\seToolBare}{0.09}
\newcommand{\nToolBare}{96}
\newcommand{\cMatched}{0.32}
\newcommand{\cMatchedse}{0.05}
\newcommand{\cMatchedassign}{0.06}
\newcommand{\cMatchedp}{<0.001}
\newcommand{\cMatchedn}{128}
\newcommand{\cMatchedlo}{+0.20}
\newcommand{\cMatchedhi}{+0.44}
\newcommand{\cOrig}{-0.01}
\newcommand{\cOrigse}{0.07}
\newcommand{\cOrigassign}{0.11}
\newcommand{\cOrigp}{0.927}
\newcommand{\cOrign}{128}
\newcommand{\cOriglo}{-0.22}
\newcommand{\cOrighi}{+0.21}
\newcommand{\cPeerTrail}{-0.28}
\newcommand{\cPeerTrailse}{0.07}
\newcommand{\cPeerTrailassign}{0.07}
\newcommand{\cPeerTrailp}{<0.001}
\newcommand{\cPeerTrailn}{128}
\newcommand{\cPeerTraillo}{-0.41}
\newcommand{\cPeerTrailhi}{-0.14}
\newcommand{\cToolTrail}{0.05}
\newcommand{\cToolTrailse}{0.06}
\newcommand{\cToolTrailassign}{0.06}
\newcommand{\cToolTrailp}{0.400}
\newcommand{\cToolTrailn}{128}
\newcommand{\cToolTraillo}{-0.07}
\newcommand{\cToolTrailhi}{+0.17}
\newcommand{\iMB}{0.32}
\newcommand{\iMBcase}{0.05}
\newcommand{\iMBagent}{0.09}
\newcommand{\iMBp}{<0.001}
\newcommand{\iMBnull}{0.17}
\newcommand{\iMBn}{128}
\newcommand{\iCS}{-0.28}
\newcommand{\iCScase}{0.07}
\newcommand{\iCSagent}{0.04}
\newcommand{\iCSp}{0.002}
\newcommand{\iCSnull}{0.18}
\newcommand{\iCSn}{128}
\newcommand{\iMS}{0.21}
\newcommand{\iMScase}{0.08}
\newcommand{\iMSagent}{0.10}
\newcommand{\iMSp}{0.005}
\newcommand{\iMSnull}{0.15}
\newcommand{\iMSn}{192}
\newcommand{\iPF}{-0.33}
\newcommand{\iPFcase}{0.07}
\newcommand{\iPFagent}{0.08}
\newcommand{\iPFp}{<0.001}
\newcommand{\iPFnull}{0.18}
\newcommand{\iPFn}{96}
\newcommand{\iPO}{-0.20}
\newcommand{\iPOcase}{0.07}
\newcommand{\iPOagent}{0.12}
\newcommand{\iPOp}{0.016}
\newcommand{\iPOnull}{0.16}
\newcommand{\iPOn}{96}
\newcommand{\bafSd}{0.16}
\newcommand{\bafP}{0.26}
\newcommand{\decSd}{0.064}
\newcommand{\decCells}{56}
\newcommand{\mrPB}{42}
\newcommand{\mrTB}{23}
\newcommand{\mrPD}{6}
\newcommand{\mrTD}{64}
\newcommand{\mrCD}{31}
\newcommand{\bmOpus}{0.32}
\newcommand{\bmOpusse}{0.05}
\newcommand{\bmOpusn}{128}
\newcommand{\bmSonnet}{0.02}
\newcommand{\bmSonnetse}{0.21}
\newcommand{\bmSonnetn}{64}
\newcommand{\bmHaiku}{0.12}
\newcommand{\bmHaikuse}{0.07}
\newcommand{\bmHaikun}{64}
\newcommand{\confCorr}{-0.82}
\newcommand{\confGuided}{7.18}
\newcommand{\confUnguided}{6.52}
\newcommand{\cfPB}{6.00}
\newcommand{\exPB}{25}
\newcommand{\cfCD}{6.84}
\newcommand{\exCD}{8}
\newcommand{\cfTB}{6.61}
\newcommand{\exTB}{28}
\newcommand{\cfPM}{6.34}
\newcommand{\exPM}{29}
\newcommand{\cascN}{96}
\newcommand{\cascBeta}{1.00}
\newcommand{\cascBetase}{0.05}
\newcommand{\cascBetaFE}{0.55}
\newcommand{\cascBetaFEse}{0.06}
\newcommand{\cascBetan}{80}
\newcommand{\cascCopy}{71}
\newcommand{\plCopy}{48}
\newcommand{\plBeta}{0.75}
\newcommand{\plBetalo}{0.57}
\newcommand{\plBetahi}{0.97}
\newcommand{\cascCopyOne}{50}
\newcommand{\cascUnan}{5}
\newcommand{\cascCopyPosTwo}{56}
\newcommand{\cascCopyPosThree}{88}
\newcommand{\cascCopyPosFive}{81}
\newcommand{\cascPosFirst}{0.19}
\newcommand{\cascPosLast}{0.12}
\newcommand{\cascEarly}{0.056}
\newcommand{\cascLate}{0.017}
\newcommand{\cascCases}{16}
\newcommand{\cascBase}{0.286}
\newcommand{\cascFold}{17}
\newcommand{\cascGuideSd}{0.029}
\newcommand{\cascFirst}{0.06}
\newcommand{\cascFirstse}{0.04}
\newcommand{\dispMin}{0.07}
\newcommand{\dispMax}{0.33}
\newcommand{\Ndecisions}{2838}
\newcommand{\Nprimary}{1014}
\newcommand{\Ncases}{16}
\newcommand{\Nagents}{15}
\newcommand{\Nmodels}{1}
\newcommand{\dispLooLo}{0.074}
\newcommand{\dispLooHi}{0.327}
\newcommand{\dispFull}{0.286}
\newcommand{\agentFloor}{0.264}
\newcommand{\xgOpus}{-0.02}
\newcommand{\xgOpusse}{0.04}
\newcommand{\xgOpusn}{192}
\newcommand{\xgSonnet}{0.20}
\newcommand{\xgSonnetse}{0.06}
\newcommand{\xgSonnetn}{32}
\newcommand{\xgHaiku}{0.21}
\newcommand{\xgHaikuse}{0.11}
\newcommand{\xgHaikun}{32}
\newcommand{\pfPrem}{0.05}
\newcommand{\pfPremp}{0.424}
\newcommand{\pfPremnull}{0.13}
\newcommand{\poPrem}{0.18}
\newcommand{\poPremp}{0.012}
\newcommand{\poPremnull}{0.15}
\newcommand{\pdPrem}{0.04}
\newcommand{\pdPremp}{0.488}
\newcommand{\pdPremnull}{0.12}
\newcommand{\plBetaFE}{-0.97}
\newcommand{\plBetaFElo}{-2.01}
\newcommand{\plBetaFEhi}{-0.13}
\newcommand{\pullToolBareEst}{0.56}
\newcommand{\cOrigNull}{0.17}
\newcommand{\cOrigRIp}{0.945}
\newcommand{\pullPMatch}{0.76}
\newcommand{\sePMatch}{0.07}
\newcommand{\nPMatch}{96}
\newcommand{\pullPFree}{0.59}
\newcommand{\sePFree}{0.11}
\newcommand{\nPFree}{48}
\newcommand{\pullPOwn}{0.72}
\newcommand{\sePOwn}{0.11}
\newcommand{\nPOwn}{48}
\newcommand{\cascCopylo}{61}
\newcommand{\cascCopyhi}{81}

% statistics added in revision, generated by analysis/revision_macros.py in the
% artifact repository from out_robustness.json and out_power.json
\newcommand{\rbSentN}{3807}
\newcommand{\rbDistinct}{93}
\newcommand{\rbDivSix}{65}
\newcommand{\rbDivTwelve}{42}
\newcommand{\rbDivFive}{27}
\newcommand{\rbMode}{48}
\newcommand{\rbAncSix}{16.7}
\newcommand{\rbAncFive}{16.7}
\newcommand{\rbAncAll}{0}
\newcommand{\rbExact}{29}
\newcommand{\rbExactN}{1108}
\newcommand{\rbExactRound}{45}
\newcommand{\rbMatchRound}{40}
\newcommand{\rbMatchPlain}{18}
\newcommand{\rbNRMS}{0.19}
\newcommand{\rbNRMSp}{0.238}
\newcommand{\rbNRMSn}{96}
\newcommand{\rbNRMB}{0.35}
\newcommand{\rbNRMBp}{0.007}
\newcommand{\rbNRMBn}{64}
\newcommand{\rbNROS}{-0.12}
\newcommand{\rbNROSp}{0.312}
\newcommand{\rbNROSn}{64}
\newcommand{\rbMagMSLo}{0.02}
\newcommand{\rbMagMSLop}{0.869}
\newcommand{\rbMagMSHi}{0.25}
\newcommand{\rbMagMSHip}{0.005}
\newcommand{\rbMagMBLo}{0.35}
\newcommand{\rbMagMBLop}{0.039}
\newcommand{\rbMagMBHi}{0.32}
\newcommand{\rbMagMBHip}{0.001}
\newcommand{\rbMagOSLo}{-0.14}
\newcommand{\rbMagOSLop}{0.578}
\newcommand{\rbMagOSHi}{0.03}
\newcommand{\rbMagOSHip}{0.780}
\newcommand{\rbLogMS}{0.17}
\newcommand{\rbLogMSp}{0.029}
\newcommand{\rbRankMS}{0.25}
\newcommand{\rbRankMSp}{0.040}
\newcommand{\rbLogMB}{0.31}
\newcommand{\rbLogMBp}{<0.001}
\newcommand{\rbRankMB}{0.55}
\newcommand{\rbRankMBp}{<0.001}
\newcommand{\rbLogOS}{0.01}
\newcommand{\rbLogOSp}{0.892}
\newcommand{\rbRankOS}{-0.01}
\newcommand{\rbRankOSp}{0.982}
\newcommand{\rbConfDistinct}{4}
\newcommand{\rbConfTop}{98}
\newcommand{\decodeSD}{0.064}
\newcommand{\decodeCells}{56}
\newcommand{\pwUnder}{3}
\newcommand{\pwTotal}{20}
\newcommand{\pwBench}{0.32}
\newcommand{\pwMdeSonNG}{0.81}
\newcommand{\pwEstSonNG}{0.07}
\newcommand{\pwLoSonNG}{-0.50}
\newcommand{\pwHiSonNG}{0.63}
\newcommand{\pwMdeHaiNG}{0.33}
\newcommand{\pwEstHaiNG}{0.07}
\newcommand{\pwLoHaiNG}{-0.16}
\newcommand{\pwHiHaiNG}{0.30}
\newcommand{\pwMdeSonG}{0.16}
\newcommand{\pwEstSonG}{0.20}
\newcommand{\pwLoSonG}{0.09}
\newcommand{\pwHiSonG}{0.31}
\newcommand{\pwMdeHaiG}{0.31}
\newcommand{\pwEstHaiG}{0.21}
\newcommand{\pwLoHaiG}{-0.01}
\newcommand{\pwHiHaiG}{0.43}
\newcommand{\pwMdeGxSon}{0.43}
\newcommand{\pwEstGxSon}{0.63}
\newcommand{\pwLoGxSon}{0.33}
\newcommand{\pwHiGxSon}{0.93}
\newcommand{\pwMdeMS}{0.22}
\newcommand{\pwEstMS}{0.21}
\newcommand{\pwLoMS}{0.05}
\newcommand{\pwHiMS}{0.36}
\newcommand{\pwMdeMB}{0.13}
\newcommand{\pwEstMB}{0.32}
\newcommand{\pwLoMB}{0.23}
\newcommand{\pwHiMB}{0.41}
\newcommand{\pwMdePool}{0.32}
\newcommand{\pwEstPool}{0.14}
\newcommand{\pwLoPool}{-0.08}
\newcommand{\pwHiPool}{0.36}
\newcommand{\blBench}{0.206}
\newcommand{\blSonMSN}{64}
\newcommand{\blSonMSEst}{-0.142}
\newcommand{\blSonMSP}{0.480}
\newcommand{\blSonMSMde}{0.622}
\newcommand{\blSonMSLo}{-0.577}
\newcommand{\blSonMSHi}{0.293}
\newcommand{\blHaiMSN}{63}
\newcommand{\blHaiMSEst}{-0.012}
\newcommand{\blHaiMSP}{0.886}
\newcommand{\blHaiMSMde}{0.219}
\newcommand{\blHaiMSLo}{-0.165}
\newcommand{\blHaiMSHi}{0.141}
\newcommand{\blSonMBN}{64}
\newcommand{\blSonMBEst}{0.025}
\newcommand{\blSonMBP}{0.892}
\newcommand{\blSonMBMde}{0.583}
\newcommand{\blSonMBLo}{-0.383}
\newcommand{\blSonMBHi}{0.433}
\newcommand{\blHaiMBN}{64}
\newcommand{\blHaiMBEst}{0.121}
\newcommand{\blHaiMBP}{0.036}
\newcommand{\blHaiMBMde}{0.186}
\newcommand{\blHaiMBLo}{-0.009}
\newcommand{\blHaiMBHi}{0.251}
\newcommand{\eqN}{128}
\newcommand{\eqEst}{-0.005}
\newcommand{\eqP}{0.973}
\newcommand{\eqMde}{0.347}
\newcommand{\eqLo}{-0.247}
\newcommand{\eqHi}{0.238}
\newcommand{\rlN}{128}
\newcommand{\rlEst}{0.001}
\newcommand{\rlP}{0.995}
\newcommand{\rlMde}{0.390}
\newcommand{\rlLo}{-0.272}
\newcommand{\rlHi}{0.274}
\newcommand{\rlPeerN}{128}
\newcommand{\rlPeerEst}{0.012}
\newcommand{\rlPeerP}{0.905}
\newcommand{\rlPeerMde}{0.297}
\newcommand{\rlPeerLo}{-0.196}
\newcommand{\rlPeerHi}{0.220}
\newcommand{\rlToolN}{128}
\newcommand{\rlToolEst}{0.006}
\newcommand{\rlToolP}{0.945}
\newcommand{\rlToolMde}{0.254}
\newcommand{\rlToolLo}{-0.171}
\newcommand{\rlToolHi}{0.184}
\newcommand{\xlMSN}{832}
\newcommand{\xlMSEst}{0.220}
\newcommand{\xlMSP}{<0.001}
\newcommand{\xlMSMde}{0.136}
\newcommand{\xlMSLo}{0.125}
\newcommand{\xlMSHi}{0.316}
\newcommand{\xlMSBench}{0.206}
\newcommand{\xlMBN}{832}
\newcommand{\xlMBEst}{0.256}
\newcommand{\xlMBP}{<0.001}
\newcommand{\xlMBMde}{0.138}
\newcommand{\xlMBLo}{0.159}
\newcommand{\xlMBHi}{0.353}
\newcommand{\xlMBBench}{0.323}
\newcommand{\solMSN}{832}
\newcommand{\solMSEst}{0.239}
\newcommand{\solMSP}{<0.001}
\newcommand{\solMSMde}{0.150}
\newcommand{\solMSLo}{0.134}
\newcommand{\solMSHi}{0.344}
\newcommand{\solMSBench}{0.206}
\newcommand{\solMBN}{832}
\newcommand{\solMBEst}{0.113}
\newcommand{\solMBP}{<0.001}
\newcommand{\solMBMde}{0.123}
\newcommand{\solMBLo}{0.028}
\newcommand{\solMBHi}{0.199}
\newcommand{\solMBBench}{0.323}
\newcommand{\rSevenPairs}{312}
\newcommand{\rSevenN}{624}
\newcommand{\rSevenEst}{+0.250}
\newcommand{\rSevenP}{<0.001}
\newcommand{\rSevenMde}{0.161}
\newcommand{\rSevenLo}{+0.134}
\newcommand{\rSevenHi}{+0.365}
\newcommand{\rSevenPeerSlope}{0.911}
\newcommand{\rSevenToolSlope}{0.661}
\newcommand{\rSevenPeerExact}{50.6}
\newcommand{\rSevenToolExact}{15.4}
\newcommand{\rSevenPeerOutside}{23.1}
\newcommand{\rSevenToolOutside}{71.8}


% left-justify all captions (the article default centers captions short enough to
% fit on one line, which leaves Figure 9 and Table 4 inconsistent with the rest);
% fonts, sizes and skips are unchanged from the style file
\makeatletter
\long\def\@makecaption#1#2{%
  \vskip\abovecaptionskip
  \noindent #1: #2\par
  \vskip\belowcaptionskip}
\makeatother

\title{Wording Rivals Peers in LLM Judicial Decision-Making}
\author{%
  Austin L. Cheng \\
  Stanford Law School \\
  Stanford, CA 94305 \\
  \texttt{austincheng98@gmail.com} \\
}
\workshoptitle{SocialAgent}
\AtEndDocument{\typeout{PAGECOUNT \thepage}}
\begin{document}
\maketitle
\raggedbottom

\begin{abstract}
Language-model agents often converge after observing others' decisions, but this pattern can reflect either
peer influence or numerical anchoring. Convergence alone cannot distinguish them. We test this in fictional
criminal sentencing with a yoked dose-response design. LLM judges see three numbers displaced from the
guideline midpoint, attributed in paired conditions to peer sentences or a statistical forecast. Separate
dose-response slopes estimate numerical pull under peer and forecast attribution, with their difference
being the peer premium. With the guideline omitted, Claude Opus 5's premium was $\iMB$ under bare blocks
and $\iMS$ under structure-matched blocks. Under identical prompts and case allocations, GPT-6 Luna and Sol
had matched premiums of $\xlMSEst$ and $\solMSEst$ (both $p<0.001$). Removing the guideline raises pull
across all three Claude models. In civil damages, the estimated peer-minus-forecast difference in mean
proportional distance from the displayed center was $+0.140$ for Opus 5 and $+0.093$ for GPT-6 Sol.
Together, these results extend the peer premium to another model lineage and the peer-versus-forecast
comparison to a separate legal domain, while showing that wording and guideline availability affect measured
influence. The 4,086 core decision records, harness, analyses, and prompts are available at
\url{https://github.com/austin-cheng98/llm-sentencing-agents}.
\end{abstract}

\section{Introduction}

Populations of language-model agents are increasingly used to study social processes alongside
experiments with human participants. Simulated populations reproduce survey responses
\citep{argyle2023out,park2024thousand}, replicate classic behavioral experiments
\citep{aher2023using,horton2023large}, and model opinion dynamics on networks
\citep{chuang2024simulating}. In these studies, agents who observe one another often move toward one
another. Researchers often interpret this convergence as conformity, drawing on studies of group
influence \citep{asch1956independence,sherif1935study} and information cascades
\citep{bikhchandani1992theory,anderson1997information}.

Convergence alone does not identify peer influence. When an agent sees a peer's decision, its prompt
contains both a social source and a number. Language models can move toward numbers presented in
context regardless of their source \citep{jones2022capturing,echterhoff2024cognitive}; they can also
move toward an interlocutor's position \citep{sharma2024sycophancy}. Both mechanisms predict
convergence. Prior studies distinguish them using discrete outcomes and varying whether a source is
present \citep{hu2026conformity,zhong2025disentangling}. We instead use a continuous outcome and vary
the displacement of displayed numbers. The resulting slope estimates numerical pull; the difference
between source-specific slopes estimates the peer premium.

We study this distinction in fictional sentencing decisions. Paired conditions use the same cases
and displayed numbers but attribute them either to other judges or to a statistical forecast. On
Claude Opus 5 with the guideline removed, the premium is $\iMB$ under bare blocks and $\iMS$ under
structure-matched blocks. Adding the routine instruction ``You are deciding the same case
independently'' reduces it to $\cOrig$. The instruction changes pull by $\iCS$; its two paraphrases
change pull by $\iPF$ and $\iPO$. \citet{weng2025conformity} raised independence rates with a
persona instruction. Here, the routine sentence reduces the premium on Opus 5 but not on the two
smaller models. Removing the guideline raises pull across the three Claude models. In a sequential
cascade, $\cascCopy\%$ of decisions exactly repeat a predecessor, compared with a $\plCopy\%$
exact-match rate when independent agents are shuffled into speaking orders.

Sentencing provides a continuous outcome, an institutional anchor that can be removed, and prior
evidence that human sentencing decisions respond to numerical anchors
\citep{englich2006playing,guthrie2001inside}. Language models have also been proposed as drafting
aids in this setting.

\section{Related work}

Conditioning a model on demographic profiles reproduces human subgroup response distributions
\citep{argyle2023out,park2024thousand}. Agent architectures with memory and reflection produce
social interactions \citep{park2023generative}, and classic experiments have been
replicated with simulated subjects \citep{aher2023using,horton2023large}. However,
\citet{santurkar2023whose} document systematic gaps between model opinion distributions and their
targets. Model outputs also change with wording: \citet{sclar2024quantifying} show that formatting
alone can change benchmark accuracy, and \citet{mizrahi2024state} argue that single-template
evaluation misstates model rankings. We ask whether a routine instruction changes the measured size
of a peer effect, rather than task accuracy.
\citet{chuang2024simulating} report that LLM agents in network simulations converge toward accurate
consensus; prompts inducing confirmation bias produce fragmentation. \citet{zhou2025pimmur} find
that collective effects can vanish or reverse once validity
conditions are enforced. Two non-social mechanisms could produce convergence: numeric anchoring
\citep{jones2022capturing,echterhoff2024cognitive} and movement toward an interlocutor's expressed
views \citep{sharma2024sycophancy}. In humans, \citet{englich2006playing} found experienced judges
shifting sentencing recommendations toward dice-roll anchors, and \citet{guthrie2001inside}
documented anchoring among United States federal magistrate judges.

\section{Experimental design}

\subsection{Task and cases}

Decisions come from Claude Opus 5, Sonnet 5 or Haiku 4.5 at default decoding settings, one draw
per design cell, and a GPT-6 Luna sample (Section~\ref{sec:across}). Repeat draws on
a byte-identical prompt, available for \decCells\ cells collected twice, differ by a within-cell
standard deviation of \decSd.

Each agent acts as a sentencing judge in a fictional jurisdiction. It returns an integer sentence
in months, a confidence rating, a primary and a secondary rationale from a fixed six-item menu, and
a short justification. Every prompt states that the cases are procedurally generated and that no
real defendant or court is involved.

We generate \Ncases\ cases by fully crossing four binary factors: offense severity, prior record,
expressed remorse, and substantial assistance to investigators. Offense type rotates across four
families, so each type co-occurs equally often with each factor level. Severity and prior record
determine an advisory range taken from the United States Sentencing Guidelines table, giving four
ranges below the statutory maxima. Remorse and assistance do not appear in the range, so they can
affect sentences only through departures from the guideline.

Before deciding, agents in the anchor arms see three numbers centred on a displaced point,
\begin{equation}
\bar a_{ic} \;=\; m_c\,(1 + \delta_{ic}),
\qquad
\delta_{ic} \;=\; D\bigl[(g_c + i) \bmod 4\bigr],
\qquad
D = (-0.30,\; -0.15,\; +0.15,\; +0.30),
\label{eq:anchor}
\end{equation}
where $m_c$ is the midpoint of case $c$'s advisory range, $g_c \in \{0,1,2,3\}$ is its design group
and $i$ indexes the agent.

Appendix~\ref{app:design} gives the full rotation in Equation~\ref{eq:anchor}. The sixteen design
points are partitioned into four groups in which every case
factor splits evenly. Each agent therefore sees every displacement exactly four times, and the four
factors split evenly within each displacement level. As a result, $\hat\pi$ carries no part of a
factor effect. The rotation has period four, so with six agents the first two displacement columns
repeat.

\subsection{Conditions and estimation}

We cross attribution with the presence of an institutional anchor. In the peer condition, the
numbers are sentences already entered in this case by other judges of the same bench. In the
forecast condition, they are the output of a regression model fitted to historical dispositions. In
the docketing condition, they are stored values attached to the file by the case-management system,
with no legal significance. The guideline-present condition states the advisory range; the
guideline-removed condition states only the statutory maximum. Each decision uses a fresh context,
preventing information carryover across prompts.

Each block originally ended with a closing sentence, but the sentences were not matched. The peer
block closed with ``You are deciding the same case independently,'' whereas the forecast block ended
with a note that the forecast is advisory. We added arms with both closing sentences removed, an arm
that gave the peer block a descriptive sentence matching the forecast block's structure, and two arms
with paraphrases of the independence sentence. The artifact repository's \texttt{experiment/prompts.txt} file
gives all nine prompts verbatim.

The outcome is the proportional deviation of the sentence from the case's guideline midpoint,
$y_{ijc} = (s_{ijc} - m_c)/m_c$, which puts cases with different ranges on a common scale. The
quantity of interest is the pull $\pi$,
\begin{equation}
y_{ijc} \;=\; \alpha \;+\; \pi\,\delta_{ic} \;+\; \gamma^{\top} F_c \;+\; \varepsilon_{ijc},
\label{eq:pull}
\end{equation}
where $F_c$ collects the four case factors. $\pi = 0$ means the agent ignores the numbers and
$\pi = 1$ means it adopts them in full. Writing $P_{j} = 1$ for the peer arm and $0$ for the
forecast arm, the peer premium $\Delta$ is the difference in slopes recovered by interacting $P_j$
with $\delta_{ic}$ in Equation~\ref{eq:pull},
\begin{equation}
y_{ijc} \;=\; \alpha \;+\; \pi\,\delta_{ic} \;+\; \beta P_{j} \;+\;
\Delta\,(\delta_{ic} \times P_{j}) \;+\; \gamma^{\top} F_c \;+\; \varepsilon_{ijc},
\label{eq:premium}
\end{equation}
so that $\Delta = \partial^2 \mathbb{E}[y] / \partial \delta \, \partial P$. A moderation model
adding a guideline indicator is reported in Appendix~\ref{app:inference}; only the arms carrying
the original closing lines were run on both sides of it, so we report the two guideline conditions
separately instead.

Paired arms share agents, cases, and displacements, so each cell $(i,c)$ contributes one decision
under each label. Under the sharp null that a decision would have been the same under either label,
$y_{ic}(P{=}1) = y_{ic}(P{=}0)$ for every cell, the labels within a cell are exchangeable;
relabelling them at random per cell generates the exact null distribution of $\hat\Delta$, from
which we calculate a two-sided $p$-value and a 95th-percentile critical value over $B = 9{,}999$ draws.
Each contrast has its own critical value because it rests on a different set of cells. The values
range from $\pdPremnull$ to $\iCSnull$ here, and each premium is evaluated against its corresponding
critical value. Because the contrast never crosses agents, between-agent heterogeneity cancels from
the null rather than inflating it. Standard errors cluster on case, agent, and displacement cell.
Appendix~\ref{app:inference} gives the relabelling scheme, the wild cluster bootstrap used for
cluster-robust $p$-values, the per-arm displacement null, the between-agent spread $\bar\sigma$, and
the yoke, factor-balance, and offense-type checks.

\subsection{Civil-damages extension}

The extension tests monetary valuations on 16 fictional bodily-injury claims in a $2^4$ factorial
crossing severity, liability clarity, claimant mitigation, and documentation completeness. Each
claim has a stipulated midpoint and an advisory range of $\pm12.5\%$, but the range is withheld;
agents may enter any amount up to the one-million-dollar policy limit. Prompts show three valuations
at 94\%, 100\%, and 106\% of a midpoint shifted by $\delta\in\{-0.30,-0.15,0.15,0.30\}$, rounded
to the nearest hundred dollars. Matched primary arms use identical cases and figures and the same
no-other-information cue; only attribution (other adjusters or a regression forecast) changes. We
measure in-band awards, proportional distance from the displayed center, awards within 2\% of
center, and exact matches. D1's initial Opus~5 panel informed outcome selection; the measures and
directions were registered before examining two additional agents (32 pairs). Preregistered D2
reused the same 320 prompts and assignments with GPT-6 Sol high; one failed response left 159 pairs.

The in-band indicator includes both displayed endpoints. Proportional distance is the absolute
deviation from the central figure divided by that figure; we also code whether the award falls
within 2\% of center and whether it matches the center exactly. Binary outcomes use two-sided exact
sign tests on discordant matched pairs. For distance, we estimate the mean paired peer-minus-tool
difference, cluster uncertainty by claim, and report the registered minimum detectable effect.

\section{Results}

\subsection{The peer premium}

On Claude Opus 5 with bare blocks, peer attribution gives a pull of $\pullPeerBare$, compared with
$\pullToolBareEst$ for the forecast,
yielding a premium of $\iMB$ (SE \iMBcase\ on case, \iMBagent\ on agent, $p\,\iMBp$, $n=\iMBn$).
The peer block has four lines, whereas the forecast block has five because naming a forecast requires
a descriptive sentence. Matching the structures reduces the premium to $\iMS$ (SE \iMScase\ on case,
$p=\iMSp$), against a design null of \iMSnull; restoring the closing sentences gives $\cOrig$.
The lower structure-matched estimate is consistent with the unmatched block structures contributing
to the bare comparison. The bare comparison on two smaller models gives $\bmHaiku$ and $\bmSonnet$;
the last four rows report the GPT-6 cross-lineage estimates (Section~\ref{sec:across};
Figure~\ref{fig:forest}).

\begin{figure}[H]
\centering
\includegraphics[width=1.0\textwidth]{../figures/fig_forest.pdf}
\caption{Every estimate compares peer attribution against a statistical forecast over
bit-identical numbers, with the advisory guideline removed throughout. Shading marks the design
null of that row's own contrast, the 95th percentile of its label-swap distribution. Band widths
vary because the contrasts have different design-specific null distributions. GPT-6 rows have no
label-swap band. Intervals are 95\% and cluster on case.}
\label{fig:forest}
\end{figure}

Across bare-block arms, premiums are $\iMB$ on Opus 5, $\bmHaiku$ on Haiku 4.5, and $\bmSonnet$ on
Sonnet 5. The ordering reverses with the guideline: Opus 5 is at $\xgOpus$, versus $\xgSonnet$ and
$\xgHaiku$ in cells containing \xgSonnetn\ decisions from one agent (Appendix~\ref{app:models}).
These cells are too small to rank the models. The independent-decision sentence reduces the Opus 5
premium but not those in the smaller-model arms, so comparisons depend on wording.

\subsection{Instruction wording}
\label{sec:sentence}

``You are deciding the same case independently'' changes pull by $\iCS$ ($p=\iCSp$, $n=\iCSn$); ``Form your own
view of what this case warrants'' changes it by $\iPF$ ($p\,\iPFp$, $n=\iPFn$); and ``Your sentence should
reflect your own judgment of the case'' changes it by $\iPO$ ($p=\iPOp$). Each change is statistically
distinguishable from its corresponding label-swap null. None names a peer, withdraws the block, or changes a displayed number. The forecast hedge leaves
pull unchanged ($\cToolTrail$, $p=\cToolTrailp$). Because a premium is a difference in pulls, these
effects account for the premium observed with both original sentences in place:
\begin{equation}
\hat\Delta_{\text{orig}} \;=\; \hat\Delta_{\text{bare}} \;+\; \hat c_{\text{peer}}
\;-\; \hat c_{\text{fcst}}
\;=\; \iMB \;+\; (\iCS) \;-\; (\cToolTrail) \;=\; \cOrig,
\label{eq:decomp}
\end{equation}
Because the paired prompts keep cases and displayed numbers fixed, this decomposition estimates the
response to wording changes within this task. It does not distinguish whether the added sentence
changes perceived independence, the social meaning assigned to peer sentences, or attention to the
anchors. The near-zero forecast-hedge estimate provides a comparison: adding a sentence did not
reliably change pull in that arm. The results locate a wording effect in the peer condition while
leaving its psychological mechanism unmeasured (Figure~\ref{fig:sentences}).

\begin{figure}[H]
\centering
\includegraphics[width=0.75\textwidth]{../figures/fig_sentences.pdf}
\caption{Pull before and after the sentence is appended, with the numbers shown to the agent
held identical throughout.}
\label{fig:sentences}
\end{figure}

Measured against the same bare forecast, the two paraphrases have different effects: one reduces the
estimated premium to $\pfPrem$, whereas the other leaves it at $\poPrem$
(Appendix~\ref{app:selfreport}). All three wordings lower peer pull, which points to the
instruction's content rather than its exact wording.

The sentence also changes justifications: peer mentions fall from \mrPB\% to \mrPD\% when it is
added, alongside the reduction in pull, although the text does not establish how agents interpreted
the instruction.

\subsection{The advisory guideline}

Removing the advisory range raises estimated pull in each of the three Claude models. The pooled
interaction is $\ambig$ (SE \ambigse, $p\,\ambigp$). By model, the estimates are
$\gxOpus$ for Opus 5, $\gxSonnet$ for Sonnet 5, and $\gxHaiku$ for Haiku 4.5
(Figure~\ref{fig:models}, Appendix~\ref{app:models}). Only arms carrying the original closing lines
were run with the guideline, so this contrast holds those lines fixed rather than spanning the
wordings in Section~\ref{sec:sentence}. Pull under a guideline is $\gpOpus$, $\gpSonnet$, and
$\gpHaiku$.
Agent fixed effects over the \feAgents\ agents in both settings leave the interaction at
\feAmbig\ (SE \feAmbigse).

\begin{figure}[H]
\centering
\includegraphics[width=1.0\textwidth]{../figures/fig_models.pdf}
\caption{Pull with 95\% intervals clustered on case. (a) Original closing sentences, with (solid)
and without (hatched) the advisory guideline. (b) Bare blocks, no guideline, collected for every model.}
\label{fig:models}
\end{figure}

Between-agent spread is $\NaGsd$ with the guideline and $\NaNsd$ without it. Leave-one-out estimates
for the no-guideline condition range from $\dispLooLo$ to $\dispLooHi$
(Appendix~\ref{app:models}). Without a guideline, mean deviation is $\NaNdev$ compared with
$\NaGdev$; displaying the three numbers shifts it to $\IntTN$.

\subsection{The docketing arm}

A docketing condition provides a source comparison with no legal significance. Presented as
case-management records, the same values move agents by $\pullClerN$ (SE \seClerN, $n=\nClerN$),
compared with $\pullToolN$ for the forecast ($\twCler$, SE \twClerse, $p\,\twClerp$). The prompt
also tells agents the values are unrelated to the merits, so the estimate combines the source label
with the instruction to disregard the numbers (Figure~\ref{fig:main}).

\begin{figure}[H]
\centering
\includegraphics[width=1.0\textwidth]{../figures/fig_main.pdf}
\caption{Small points are decisions; large points are cell means; the diagonal marks full
adoption. (a) Opus 5, original blocks, guideline present (solid) or removed (hatched); (b) Opus 5,
bare blocks; (c) GPT-6 Luna and Sol pooled within each bare-block arm, with model-specific
intercepts and a common slope; (d) selected pull estimates with 95\% intervals clustered on case.}
\label{fig:main}
\end{figure}

Case-factor estimates show that sentence length varies with case characteristics even when no
numbers are shown (Appendix~\ref{app:factors}). Severity and prior record have larger effects than
assistance and remorse. With a guideline, most sentences fall within its range; without it,
proportionality is less commonly named and deterrence is more common.

\subsection{Confidence and justifications}
\label{sec:saidleast}

Each decision has a confidence rating and a written justification. Mean confidence is negatively
correlated with pull across the six unguided anchor arms of Claude Opus 5 ($r=\confCorr$;
Figure~\ref{fig:confidence}, Appendix~\ref{app:selfreport}). The bare peer arm has the highest
pull and lowest confidence, at \cfPB; the docketing arm has the lowest pull and highest confidence,
at \cfCD. Claude Sonnet 5 and Haiku 4.5, and GPT-6 Luna and Sol, are plotted separately for their
available shared arms; only Opus 5 enters the fit. Panel~\ref{fig:confidence}b averages exact adoption
by provider across the three arms shared by all five models. The correlation is descriptive and based
on six arm means; confidence does not identify decisions unaffected by the manipulation.
\label{sec:confidence}

Written justifications differ by instruction even when estimated pull is similar. With the guideline
removed and each block carrying its own closing sentence, peer and forecast pull are $\pullPeerN$ and
$\pullToolN$ (difference $\cOrig$, $p=\cOrigRIp$). The peer arm names displayed numbers in
\menPN\% of justifications, compared with \menTN\% in the forecast arm. Without the peer closing
line, peer mentions occur in \mrPB\% of justifications; with the line, they occur in \menPN\%.
These differences are consistent with the instruction affecting what agents report, while the pull
estimates remain similar.

Across decisions, pooling every unguided anchor arm, justifications that
name the displayed numbers accompany a pull of \pullMention\, compared with \pullSilent\ for
justifications that do not. This comparison is descriptive because it separates decisions rather
than randomized conditions. This association echoes \citet{turpin2023language}'s finding that model
explanations may not reveal what influenced an output, though that study examined single-agent
reasoning.

\subsection{The cascade}
\label{sec:cascade}

Earlier arms used experimenter-selected numbers. The cascade measures copying under sequential
observation: six agents decide the same case one after another, and each sees the sentences produced
by earlier agents. Among the $\cascBetan$ decisions with a visible predecessor, $\cascCopy\%$
reproduce one predecessor's number exactly (binomial interval: $\cascCopylo$ to $\cascCopyhi$\%).
In $\cascUnan$ of sixteen cases, the cascade ends unanimously. The comparison floor is $\plCopy\%$,
obtained by shuffling independent agents into pseudo-speaking orders in which influence is
impossible. Between-agent spread falls from $\cascEarly$ across the first two speakers to $\cascLate$
across the final pair. Mean deviation is $\cascPosFirst$ at the first speaking position and
$\cascPosLast$ at the last. The first speaker receives no peer block and matches the no-peer arm
within $\cascFirst$
(SE \cascFirstse), so the cascade compares conditions within one arm (Figure~\ref{fig:cascade}).

\begin{figure}[H]
\centering
\includegraphics[width=0.72\textwidth]{../figures/fig_cascade.pdf}
\caption{(a) Each line is one case, centred on its own mean, with agents plotted by speaking
position; the heavy line is mean absolute deviation. (b) Between-agent spread on the same
cases, for independent agents and for the first and last pairs of a cascade.}
\label{fig:cascade}
\end{figure}

Both cascade statistics have a comparison floor, which determines the primary estimator. With case
fixed effects, the observed slope is $\cascBetaFE$ (SE \cascBetaFEse), compared with a placebo floor
of $\plBetaFE$. The plain slope is close to its own floor, so we use the fixed-effect version as
primary.
Appendix~\ref{app:cascade} gives both floors and the per-seat copying rates.

\subsection{Round numbers}
\label{sec:round}

Sentence lengths cluster at round values. Over \rbSentN\ decisions agents use \rbDistinct\ distinct values;
\rbDivSix\% are divisible by six months and \rbDivTwelve\% by twelve. A decision can therefore
match a displayed number because of this clustering rather than copying. The exact-match
rate is \rbMatchRound\% when at least one displayed number is round, compared with \rbMatchPlain\%
when none is. We therefore compare copying with a shuffled floor that preserves this preference.
The peer premium is unaffected because both arms see identical numbers. Only \rbAncSix\% of the
displayed numbers are divisible by six, approximately the chance rate.

Among records in which no displayed number is divisible by six, the bare premium is $\rbNRMB$
($p=\rbNRMBp$, $n=\rbNRMBn$), compared with $\iMB$ in the full sample. The structure-matched
premium is $\rbNRMS$ ($p=\rbNRMSp$, $n=\rbNRMSn$), compared with $\iMS$. The point estimates are
similar, but the matched estimate is not statistically significant in this half-sample. The log-ratio
and within-case-rank estimates remain positive for both the bare and structure-matched premiums
(Appendix~\ref{app:robust}). Confidence uses only \rbConfDistinct\ of its ten values, with
\rbConfTop\% of ratings at six or seven; Section~\ref{sec:confidence} reports arm-mean patterns,
not a calibrated measure.

\subsection{Equivalence and reliability}
\label{sec:equiv}

The structure-matched premium could reflect attribution or perceived informativeness: judges may
view peer sentences as drawing on more case information than historical forecasts. A preregistered
pair gave both blocks the same statement: ``These figures were produced from the same case file you
have been given, and from no other information about this defendant or this offense.'' Headers and
row labels were the only remaining contrasts; prompts and analyses were fixed before the $\eqN$
decisions.

The premium is $\eqEst$ ($p=\eqP$, $n=\eqN$; 95\% CI $[\eqLo,\eqHi]$). At 80\% power and a
two-sided 5\% test, the MDE is 2.80 times the clustered SE. Here it is $\eqMde$, above the $\iMS$
benchmark, and the interval contains that value. The registered stopping rule therefore leaves the
estimate imprecise; the arm was not extended after outcome inspection.

The cue also changes number use: exact reproduction falls by roughly two thirds, out-of-range
sentences rise to $80$--$84\%$, and dispersion increases (Appendix~\ref{app:equiv}). Agents say the
figures add no case information because they derive from the same file. The result combines source
attribution with the no-other-information instruction and does not identify the source effect when
agents use the numbers more.

A separate GPT-6 Luna follow-up used neutral access wording, describing both sources as having
reviewed the case file without saying they used no other information. Before inspecting outcomes, we
fixed a factor-balanced sample of 12 cases across 26 contexts ($\rSevenPairs$ matched pairs). The
premium is $\rSevenEst$ ($p\rSevenP$, 95\% CI $[\rSevenLo,\rSevenHi]$; MDE $\rSevenMde$); its
interval excludes zero but contains the $\blBench$ benchmark. Number use was higher under peer
attribution: slopes $\rSevenPeerSlope$ vs. $\rSevenToolSlope$, exact matches $\rSevenPeerExact\%$ vs.
$\rSevenToolExact\%$, and out-of-range sentences $\rSevenPeerOutside\%$ vs. $\rSevenToolOutside\%$.
This distinct prompt comparison does not resolve the registered no-other-information contrast.

A second registered pair added the same reliability claim to both blocks. Its premium is $\rlEst$
($p=\rlP$, $n=\rlN$, MDE $\rlMde$); adding the claim changes peer-arm pull by $\rlPeerEst$
($p=\rlPeerP$) and forecast-arm pull by $\rlToolEst$ ($p=\rlToolP$). All four estimates are near
zero and underpowered against $\iMS$; they do not establish whether the claim restores the premium.

\subsection{Other models}
\label{sec:across}

The structure-matched premiums on Sonnet~5 and Haiku~4.5 are imprecise against the Opus benchmark
($\blSonMSEst$, MDE $\blSonMSMde$; $\blHaiMSEst$, MDE $\blHaiMSMde$). Haiku's bare premium is
$\blHaiMBEst$ ($p=\blHaiMBP$); these smaller-model cells do not support reliable rankings.

GPT-6 Luna and Sol each contribute 1,248 decisions across 26 contexts, using identical prompts,
cases, numbers, and assignments (Table~\ref{tab:gpt-cross}). Their matched premiums are
$\xlMSEst$ and $\solMSEst$; bare premiums are $\xlMBEst$ and $\solMBEst$ (all $p<0.001$).
The MDEs are below the corresponding Opus benchmarks. The matched intervals include those benchmarks,
while Sol's bare interval is below the Opus estimate. Luna used extra-high and Sol high reasoning;
their estimate difference is descriptive and does not isolate model identity.

\subsection{Civil-damages extension}

In the registered Opus~5 confirmation, 14/32 peer awards fell within the range shown in the prompt,
compared with 26/32 forecast awards (exact sign test, $p=0.0042$). Peer awards were also less likely
to fall within 2\% of the center (10/32 vs. 24/32; $p=0.00052$).

Exact matches were less frequent under peer attribution (2/32 vs. 23/32; $p<0.000001$). Mean
proportional distance from the center was $0.217$ for peer awards and $0.077$ for forecast awards
(difference $+0.140$, 95\% CI $[0.032,0.248]$), just below the registered MDE of $0.142$.

The full D1 panel is descriptive because its first eight agents informed outcome selection. Across
that panel, peer awards were less close to the center on all three measures;
mean distance was $0.392$, compared with $0.156$ for forecast awards (difference $+0.236$, 95\% CI
$[0.145,0.328]$).

Preregistered D2 used the same 320 prompts and assignments with GPT-6 Sol high. One failed response
left 159 matched pairs. In-band rates were similar (90/159 vs. 88/159; $p=0.871$), as were rates
within 2\% of center and exact matches (Figure~\ref{fig:damages-civil}; $p=1.000$ and $p=0.200$).
Mean distance remained higher under peer attribution: $0.274$ vs. $0.182$ (difference $+0.093$, 95\%
CI $[0.011,0.174]$). This estimate was positive but below the registered MDE of $0.107$. Thus, D2
did not reproduce the in-band contrast, although its mean-distance estimate pointed in the same
direction as Opus~5 (Figure~\ref{fig:damages-civil}). Appendix~\ref{app:damages} reports full results.

\begin{figure}[H]
\centering
\includegraphics[width=1.0\textwidth]{../figures/fig_damages_civil.pdf}
\caption{Civil-damages results. Top: discordant matched pairs as a share of all pairs, with
peer-only awards extending right (blue) and forecast-only left (red). Bottom: peer-minus-forecast
distance from center with 95\% claim-clustered $t$ intervals; open diamonds mark minimum detectable
effects. Samples: Opus~5 D1 confirmation ($n=32$), full D1 panel ($n=160$, descriptive), and GPT-6
Sol D2 ($n=159$).}
\label{fig:damages-civil}
\end{figure}

\section{Discussion}

The central estimand is the response to a displayed number when its source attribution changes,
holding the case and number fixed. For Opus~5, matching block structure reduces the bare-block
premium from $\iMB$ to $\iMS$. The matched comparison removes differences in length and closing
sentences, making attribution a more direct interpretation of the contrast. It still measures a
response to source labels in this prompt, rather than identifying conformity as the agent's
psychological mechanism. The lower matched estimate also shows why the wording of both alternatives
is part of the estimand.

Other prompt changes locate further sources of sensitivity. The independent-judgment sentence and
its two paraphrases reduce peer pull, while the forecast hedge leaves pull nearly unchanged. Removing
the advisory guideline raises pull across the three Claude models (Figure~\ref{fig:models}b),
indicating that the institutional reference point changes how agents use the displayed values.
Docketing labels produce little pull when agents are instructed to disregard those values. In the
cascade, exact reproduction ($\cascCopy\%$) exceeds the shuffled baseline ($\plCopy\%$), while
between-agent spread falls across speaking positions. Together, these results distinguish responses
to source, wording, institutional anchors, and sequential exposure; no single convergence measure
captures all four. GPT-6 Luna and Sol also show positive premiums, extending the comparison to
another model lineage without establishing equivalence across models.

The civil-damages extension tests the same source contrast on fictional compensation awards. In the
registered Opus~5 confirmation, 14/32 peer awards, compared with 26/32 forecast awards, fell within
the displayed range. GPT-6 Sol rates were similar (90/159 vs. 88/159), so the in-range contrast did
not replicate. Mean proportional distance from the displayed center was higher under peer
attribution in both runs (+0.140 for Opus~5; +0.093 for GPT-6 Sol). The Sol estimate was positive
(95\% CI [0.011, 0.174]) but below the registered MDE (0.107). Thus, the in-range difference
appears only in Opus~5, while mean distance is higher under peer attribution in both runs. This
extension applies the design to a second legal outcome, but does not show a uniform pattern across
models and measures.

For legal simulations, these findings make source wording and institutional anchors measurable
features of the evaluation. A simulated peer effect should be reported separately from accuracy on
held-out decisions. The same distinction applies when simulated
agents are used to study social influence. Matching the displayed information and textual structure,
then varying source and context, helps identify which prompt component changes the measured response.

\section{Limitations}

The sentencing cases and civil claims are fictional, and the study has neither a human baseline nor
an external measure of decision accuracy. Its primary premium estimates how outputs respond to
attribution under controlled prompts. They do not measure whether agents use legally relevant facts.
The structure-matched estimate is also
specific to the wording and information statement tested here. Other source descriptions or
institutional settings could produce different responses. The no-other-information wording changes
number use as well as attribution, and the neutral-access follow-up uses different wording; neither
comparison isolates a general source effect across all information conditions.

Precision varies across arms. Equivalence and reliability comparisons remain underpowered, and
smaller-model estimates are too imprecise for dependable rankings. The Claude models share a lineage;
GPT-6 is a separate family, with Luna and Sol evaluated under different reasoning settings. Each
prompt cell uses one draw, so the estimates do not characterize run-to-run variability. In the
civil-damages study, the full Opus~5 D1 panel (160 pairs) is descriptive because its first eight
agents informed outcome selection; the registered confirmation uses 32 pairs. GPT-6 Sol reused the
320 prompts and assignments, with one failed response leaving 159 pairs. Its in-range rates were
similar (90/159 peer vs. 88/159 forecast), but its mean proportional-distance difference was
positive (+0.093, 95\% CI [0.011, 0.174]) and below the registered MDE (0.107). The interval
excludes zero, although the point estimate is smaller than the preregistered detection threshold.
The cascade also uses a single six-agent sequence per case, which limits inference about other group
sizes, speaking orders, and exposure patterns. These design and precision limits bound the claims to
the studied prompts, models, and fictional cases.

\section{Conclusion}

Across fictional sentencing decisions, measured peer premiums vary with source attribution, prompt
wording, and guideline availability. Structure matching clarifies the source comparison, and positive
GPT-6 estimates extend it to another model lineage. The civil-damages extension applies the design to
fictional compensation awards. The in-range difference appears in the Opus~5 confirmation but not
GPT-6 Sol; mean distance is higher under peer attribution in both. These results do not establish a
uniform effect across measures and models. The experiments measure prompt sensitivity. Evaluations of legal AI should
distinguish responses to recommendation sources from accuracy against case outcomes.

\bibliographystyle{plainnat}
\bibliography{refs}

\appendix
\raggedbottom
\makeatletter
\renewcommand{\theHfigure}{appendix.\thefigure}
\renewcommand{\theHtable}{appendix.\thetable}
\makeatother
\section{Ethical considerations}

We simulate criminal sentencing, a setting where decisions affect defendants' liberty. Every case is
procedurally generated and fictional, and every prompt says so. This study does not evaluate language
models as sentencing tools. With the guideline removed, model output responds to the arbitrary
forecast displacement ($\pullToolN$). The sentences are measurements of model behavior on synthetic
cases, not judgments about real defendants.

\section{The memory experiment}
\label{app:memory}

Deployed agent architectures more often show an agent a summary of itself than a summary of
its neighbours, so we ran a second experiment on the same cases. Agents in the
own-record arm receive a profile of their past decisions, refreshed every four
cases, giving their mean sentence, their average departure from the advisory range, how their
sentences differed by remorse and by assistance, and their most frequent rationale. Agents in
the no-memory arm decide the identical cases in the identical order with a fresh
context each time, and four baseline cases return at the end in paraphrased form, which gives
a within-agent test that no change in case composition can produce.

Without memory, the estimated trend across sixteen cases is $\driftNo$ (SE \driftNose). With an
own-record summary, the estimated trend is $\driftOwn$ (SE \driftOwnse), which is not distinguishable
from zero. The paired difference between the arms on the same case at the same position is
$\armgap$ (SE \armgapse). On returning cases, the estimated shifts are $\reexNo$ without memory
and $\reexOwn$ with memory.

We tested own-record summaries because they provide numerical information similar to peer-supplied
sentences. Each arm contains three agents and sixty decisions, limiting precision for smaller effects.

\section{What was collected}
\label{app:census}

The original arms were drawn on 17--18 August 2026 at default decoding settings, one draw per
cell, using only the prompts in the artifact repository's \texttt{experiment/prompts.txt} file; staged collection and later agent additions
explain the unequal arm sizes.

The registered cross-model arms were collected from 30 September to 1 October through the released
harness, which creates a fresh context for each prompt and inherits decoding settings. Parse failures
were recorded without redraws; one Haiku structure-matched cell failed, leaving \blHaiMSN\ rather than
\blHaiMBN\ decisions. The GPT-6 Luna sample used 26 fresh contexts across three arms and 16 cases
(1,248 decisions). It used isolated contexts with model \texttt{gpt-6-luna} and extra-high
reasoning; all replies parsed, and no cells were redrawn.

For the sentencing study, a GPT-6 Sol run was collected on 4 October with the same 1,248 cells, prompt text, cases, displayed
numbers, and displacement assignments. It used model \texttt{gpt-6-sol} at high reasoning effort.
Six first attempts returned no model output and were retried once under the protocol; all final
records parsed, and no tools were called. The Luna run used extra-high reasoning, so that setting
differs between the two model samples.

\section{Design details}
\label{app:design}

Crossing offense levels 16 and 24 with criminal history categories I and IV yields the four
advisory ranges used throughout: 21--27, 33--41, 51--63, and 77--96 months (Figure~\ref{fig:design}).

\begin{figure}[H]
\centering
\includegraphics[width=0.80\textwidth]{../figures/fig_design.pdf}
\caption{(a) Anchor values at four displacements around the advisory midpoint. (b) Identical
numbers under peer and forecast labels; their slope difference is the peer premium.}
\label{fig:design}
\end{figure}

Write $F_c \in \{0,1\}^4$ for case $c$'s factor vector and partition the sixteen design points
into four groups $g_1,\dots,g_4$ subject to
\begin{equation}
\sum_{c \in g_k} F_{cf} \;=\; \tfrac{1}{2}\,|g_k| \qquad \text{for every group } k
\text{ and every factor } f,
\label{eq:balance}
\end{equation}
which a backtracking search satisfies exactly at $|g_k| = 4$. Agent $i$ then receives
\begin{equation}
\delta_{ic} \;=\; d_{\,(k(c) + i) \bmod 4}, \qquad
d \;=\; (-0.30,\,-0.15,\,+0.15,\,+0.30),
\label{eq:alloc}
\end{equation}
where $k(c)$ indexes the group containing $c$. Every agent sees each displacement exactly four
times, and \eqref{eq:balance} makes the four factors split evenly within each displacement level.

\section{Inference}
\label{app:inference}

Drawing $\zeta_{ic} \sim \mathrm{Unif}\{0,1\}$ independently per cell and relabelling
$P^{(b)}_{ic} = \zeta_{ic} P_{ic} + (1-\zeta_{ic})(1 - P_{ic})$ generates the exact null
distribution of $\hat\Delta$, from which
\begin{equation}
p \;=\; \frac{1 + \#\{b : |\hat\Delta^{(b)}| \ge |\hat\Delta|\}}{1 + B},
\qquad
q_{.95} \;=\; \text{95th percentile of } |\hat\Delta^{(b)}|,
\label{eq:ri}
\end{equation}
with $B = 9{,}999$.

We quote all three clusterings for the premium contrasts and the case-clustered figure for per-arm
pull. Cluster-robust $p$-values use a wild cluster bootstrap with Webb six-point weights
\citep{webb2023reworking,cameron2008bootstrap,roodman2019fast}, which we use because
agent-clustered inference rests on six clusters. Per-arm pull uses a second randomization null,
permuting the displacement within each agent across the design's own factor-balanced case groups;
\eqref{eq:ri} describes the attribution-label null used for premiums. Population agreement is
summarised by the between-agent spread on a common case,
\begin{equation}
\bar{\sigma} \;=\; \frac{1}{|\mathcal{C}|}\sum_{c \in \mathcal{C}}
\operatorname{sd}_{i}\!\left( s_{ic} / m_c \right),
\label{eq:spread}
\end{equation}
the mean over cases of the standard deviation across agents of the sentence expressed in units of
the case midpoint.

Across \yokedCells\ cells in which the same agent saw the same case under more than one
attribution, the displayed numbers were identical in every paired cell. The four
case factors remain balanced across displacement levels, with the largest departure from an even
split equal to \worstImbalance. Adding offense-type fixed effects leaves the guideline-present
premium at $\rbGType$ (SE \rbGTypese).

Adding a guideline indicator $G$ with $\delta \times G$, $P \times G$ and $\delta \times P \times
G$ terms gives a moderation model whose triple interaction is $\triple$ (SE \triplese,
$p=\triplep$). Only arms with the original closing lines were run under both guideline conditions,
and their peer-premium estimates are near zero; we therefore report the conditions separately.

\section{Pull by arm}
\label{app:cells}

Arm-level pulls include all arm agents; paired premiums use shared cells (Table~\ref{tab:cells}).

\begin{table}[H]
\centering
\caption{Opus 5 pull by block and source. Case-clustered SEs; randomization $p$-values. No-number rows show mean deviation.}
\label{tab:cells}
\footnotesize
\renewcommand{\arraystretch}{0.86}
\begin{tabular}{llrrrr}
\toprule
Numbers attributed to & Block variant & $n$ & Pull $\hat{\pi}$ & SE & $p$ \\
\midrule
\multicolumn{6}{l}{\emph{Guideline removed}} \\
Other judges of this bench & none & \nPeerBare & $\pullPeerBare$ & \sePeerBare & --- \\
Other judges of this bench & structure-matched & \nPMatch & $\pullPMatch$ & \sePMatch & --- \\
Other judges of this bench & paraphrase one & \nPFree & $\pullPFree$ & \sePFree & --- \\
Other judges of this bench & paraphrase two & \nPOwn & $\pullPOwn$ & \sePOwn & --- \\
Statistical forecast & none & \nToolBare & $\pullToolBare$ & \seToolBare & --- \\
Other judges of this bench & ``decide independently'' & \nPeerN & $\pullPeerN$ & \sePeerN & \riPeerNp \\
Statistical forecast & ``advisory only'' & \nToolN & $\pullToolN$ & \seToolN & \riToolNp \\
Docketing artifact & ``no legal significance'' & \nClerN & $\pullClerN$ & \seClerN & \riClerNp \\
\addlinespace[1pt]
\multicolumn{6}{l}{\emph{Guideline present}} \\
Other judges of this bench & ``decide independently'' & \nPeerG & $\pullPeerG$ & \sePeerG & \riPeerGp \\
Statistical forecast & ``advisory only'' & \nToolG & $\pullToolG$ & \seToolG & \riToolGp \\
\midrule
\multicolumn{2}{l}{No numbers shown, guideline present} & \NaGn &
  \multicolumn{3}{l}{mean deviation $\NaGdev$} \\
\multicolumn{2}{l}{No numbers shown, guideline removed} & \NaNn &
  \multicolumn{3}{l}{mean deviation $\NaNdev$} \\
\bottomrule
\end{tabular}
\end{table}

\section{Cross-model estimates}
\label{app:models}

With a guideline and no numbers, between-agent spread is \NaGsd; without the guideline it is
\dispFull. Leave-one-out estimates for the no-guideline condition range from \dispLooLo\ to
\dispLooHi, so the estimated magnitude varies with the four-agent sample.
Guideline-present premiums per model are $\xgOpus$ (SE \xgOpusse\ on case, $n=\xgOpusn$) on Opus~5,
$\xgSonnet$ (SE \xgSonnetse) on Sonnet~5 and $\xgHaiku$ (SE \xgHaikuse) on Haiku~4.5, the latter two
over \xgSonnetn\ decisions from one agent each.

We ran both displaced-anchor arms on three models of the same family, with and without the
advisory guideline. With the guideline present,
pull increases as model size decreases, from near zero on Opus 5 to roughly two-thirds on Haiku 4.5.
Without the guideline, the spread in pull across the three models is smaller, although Opus 5 remains
below full adoption.

The estimated peer premiums are $\bmOpus$ on Opus 5, $\bmHaiku$ on Haiku 4.5, and $\bmSonnet$ on
Sonnet 5. The smaller-model estimates use $\bmSonnetn$ decisions each and are too imprecise to
separate. We make no cross-model scaling claim from these three estimates.

The GPT-6 Luna and Sol premiums use the same cases and prompts. Their full estimates are reported in
Table~\ref{tab:gpt-cross}; the comparisons across models are descriptive because the reasoning
settings differ.

\begin{table}[H]
\centering
\caption{GPT-6 peer premiums relative to the bare forecast. Each contrast uses 416 observations
per arm and case-clustered 95\% intervals. The two runs used identical prompts and allocations.}
\label{tab:gpt-cross}
\scriptsize
\setlength{\tabcolsep}{3.5pt}
\renewcommand{\arraystretch}{0.95}
\begin{tabular}{lllrcrr}
\toprule
Model & Effort & Contrast & Premium [95\% CI] & $p$ & MDE & $n$ \\
\midrule
GPT-6 Luna & xhigh & Structure-matched & $\xlMSEst$ [$\xlMSLo$, $\xlMSHi$] & $\xlMSP$ & $\xlMSMde$ & \xlMSN \\
GPT-6 Luna & xhigh & Bare & $\xlMBEst$ [$\xlMBLo$, $\xlMBHi$] & $\xlMBP$ & $\xlMBMde$ & \xlMBN \\
GPT-6 Sol & high & Structure-matched & $\solMSEst$ [$\solMSLo$, $\solMSHi$] & $\solMSP$ & $\solMSMde$ & \solMSN \\
GPT-6 Sol & high & Bare & $\solMBEst$ [$\solMBLo$, $\solMBHi$] & $\solMBP$ & $\solMBMde$ & \solMBN \\
\bottomrule
\end{tabular}
\end{table}

\section{Confidence and exact matching}
\label{app:selfreport}

Confidence uses \rbConfDistinct\ of its ten values, and
\rbConfTop\% of ratings are six or seven. We treat confidence as an uncalibrated self-report and
describe it across arm means. Against the same bare forecast, the first independence paraphrase yields
$\pfPrem$ ($p=\pfPremp$), while the second yields $\poPrem$ ($p=\poPremp$)
(Section~\ref{sec:confidence}; Figure~\ref{fig:confidence}).

\begin{figure}[H]
\centering
\includegraphics[width=0.80\textwidth]{../figures/fig_confidence.pdf}
\caption{(a) Mean confidence against pull for each model-arm cell; the dashed line and correlation
($r=\confCorr$) use only the six Opus 5 arm means (arm $n=16$--416). (b) Equal-weighted mean, by
provider, of model-level shares matching one of the three displayed numbers, for the three arms common
to all five models.}
\label{fig:confidence}
\end{figure}

\section{The cascade floors}
\label{app:cascade}

Shuffling independent agents into pseudo-speaking orders yields a \plCopy\% exact-match rate and a
slope of $\plBeta$ ($[\plBetalo, \plBetahi]$). Round-number clustering and case-level variation can
produce these matches without peer influence. With case fixed effects, the same placebo returns
$\plBetaFE$ ($[\plBetaFElo, \plBetaFEhi]$), consistent with the negative sign
\citet{angrist2014perils} predicts when case dummies place predecessor outcomes on both sides of a
within-case comparison. In the cascade, the fixed-effect slope is $\cascBetaFE$ (SE \cascBetaFEse),
compared with the plain slope $\cascBeta$ (SE \cascBetase). The fixed-effect estimate is farther
from its placebo floor; the plain estimate is close to $\plBeta$. The placebo uses at most three
predecessors, compared with five in the cascade, so it does not provide a floor matched to the later
seats. The reflection problem \citep{manski1993identification} does not arise because each agent's
predecessors are fixed before it decides.

The first four agents speak in an order randomised per case; the last two were added afterwards in
fixed positions, so position effects are identified over the first four. Exact reproduction rates
at the second, third and fourth positions are \cascCopyOne\%, \cascCopyPosTwo\% and
\cascCopyPosThree\%, respectively, with sixteen decisions at each position. In the docketing arm,
which shows three numbers in the same layout, the exact-match rate is \exCD\%. The final-pair spread
is $\cascLate$, compared with $\cascGuideSd$ for the independent population with a guideline and no
peer information. The estimates use different agent sets, so we do not calculate a fold-change.

\section{Robustness of the premium}
\label{app:robust}

\subsection{Design checks}
Across \yokedCells\ cells in which the same agent saw the same case under more than one
attribution, the numbers shown differed in \yokedMismatch\ cells, so the yoke is exact. The four
case factors remain balanced across displacement levels, with the largest departure from an even
split equal to \worstImbalance. Adding offense-type fixed effects leaves the guideline-present
premium at $\rbGType$ (SE \rbGTypese).

\subsection{Sentence values and displayed numbers}
Over \rbSentN\ decisions agents use \rbDistinct\ distinct sentence values, \rbDivSix\% divisible by
six months, \rbDivTwelve\% by twelve and \rbDivFive\% by five, with \rbMode\ months the single most
common answer. Of the numbers shown to agents, \rbAncSix\% are divisible by six and \rbAncFive\% by
five, and no record shows three round numbers at once. Decisions reproduce a displayed number in
\rbExact\% of cases (\rbExactN\ decisions), of which \rbExactRound\% are divisible by six; matching
runs \rbMatchRound\% where some displayed number is round against \rbMatchPlain\% where none is.
The original-sentence comparison is not statistically distinguishable from zero with round numbers
excluded, $\rbNROS$ ($p=\rbNROSp$, $n=\rbNROSn$). Exact matching is more common when at least one
displayed number is round than when none is. We therefore compare copying with a shuffled floor.
Because both attribution arms show the same numbers, round-number matching does not change the
estimated premium.

\subsection{Premium by displacement magnitude}
We estimate the premium at each displacement. The bare premium is flat, $\rbMagMBLo$ at
$|\delta|=0.15$ and $\rbMagMBHi$ at $|\delta|=0.30$; the structure-matched premium rises from
$\rbMagMSLo$ ($p=\rbMagMSLop$) to $\rbMagMSHi$ ($p=\rbMagMSHip$). It is larger at the far than at
the near displacement. A single slope averages these displacement-specific estimates.
The original-sentence comparison is null at both magnitudes: $\rbMagOSLo$
($p=\rbMagOSLop$) and $\rbMagOSHi$ ($p=\rbMagOSHip$).

\subsection{Alternative outcome scales}
On the log-ratio scale, the matched and bare premiums are $\rbLogMS$ ($p=\rbLogMSp$) and
$\rbLogMB$ ($p\rbLogMBp$). On the within-case midrank scale, they are $\rbRankMS$
($p=\rbRankMSp$) and $\rbRankMB$ ($p\rbRankMBp$). The original-closing comparison remains
null on both scales ($\rbLogOS$, $p=\rbLogOSp$; $\rbRankOS$, $p=\rbRankOSp$). Since ranks and
ratios use different scales, we compare signs and $p$-values, not magnitudes.

\section{Behaviour under the equivalence sentence}
\label{app:equiv}

The dispersion and copying shifts in Section~\ref{sec:equiv} are post-hoc diagnostics, not
pre-registered. On Opus~5, outcome standard deviations are $0.235$ and $0.242$ in the
structure-matched and bare pairs, versus $0.336$ and $0.329$ with the equivalence and reliability
sentences. Exact displayed-number matches are $29\%$, $25\%$, and $28\%$ in the earlier arms,
versus $8\%$, $12\%$, $6\%$, and $8\%$ in the four later arms. Sentences outside the span of the
three numbers rise from $31$--$57\%$ to $80$--$84\%$.

In both peer and forecast arms, agents independently said the numbers added no evidence beyond the
case file. One wrote that they ``derive from the identical case file and add no new information
about the defendant, so I gave them no weight as evidence.'' Responses were recorded verbatim and
not re-drawn.

\section{Case factors}
\label{app:factors}

Without numbers, severity and prior record increase sentence length by \NaNsev\% and \NaNprior\%
(Figure~\ref{fig:factors}). With a guideline, \rangeInside\% of decisions fall within its range,
\rangeBelow\% below, and \rangeAbove\% above; departures are almost entirely downward.
Proportionality is the most common stated rationale with a guideline (\RatGtopv\%); deterrence is
the most common without it (\RatNtopv\%). The factor analysis estimates larger effects for severity and prior
record (\facSev\%, \facPrior\%) than for assistance and remorse ($\facCoop$\%, \facRem\%).

\begin{figure}[H]
\centering
\includegraphics[width=0.88\textwidth]{../figures/fig_factors.pdf}
\caption{Case-factor effects on sentence length by guideline condition.}
\label{fig:factors}
\end{figure}

\section{What each arm can exclude}
\label{sec:power}

Of \pwTotal\ estimates, \pwUnder\ miss the criterion; all misses are in smaller-model cells
(Table~\ref{tab:power}).

\begin{table}[H]
\centering
\caption{Minimum detectable effect (MDE) is $2.80\times\mathrm{SE}$ for 80\% power at two-sided
5\%. ``Powered'' indicates whether each estimate meets its Claude benchmark.}
\label{tab:power}
\footnotesize
\renewcommand{\arraystretch}{0.9}
\begin{tabular}{lrrrrc}
\toprule
Estimate & $n$ & Est. & SE & MDE & Powered \\
\midrule
Structure-matched premium, Opus 5 & 192 & $+0.206$ & 0.079 & 0.222 & yes \\
Bare-block premium, Opus 5 & 128 & $+0.323$ & 0.047 & 0.132 & yes \\
Cost of the peer closing sentence & 128 & $-0.279$ & 0.074 & 0.208 & yes \\
Cost of paraphrase one & 96 & $-0.326$ & 0.066 & 0.184 & yes \\
Cost of paraphrase two & 96 & $-0.198$ & 0.070 & 0.196 & yes \\
Premium with both original closings & 128 & $-0.006$ & 0.068 & 0.191 & yes \\
Premium, Opus 5, guideline present & 192 & $-0.018$ & 0.041 & 0.116 & yes \\
Premium, Opus 5, guideline removed & 128 & $-0.006$ & 0.068 & 0.191 & yes \\
Premium, Sonnet 5, guideline present & 32 & $+0.204$ & 0.056 & 0.156 & yes \\
Premium, Sonnet 5, guideline removed & 32 & $+0.065$ & 0.290 & 0.812 & \textbf{no} \\
Premium, Haiku 4.5, guideline present & 32 & $+0.211$ & 0.112 & 0.315 & yes \\
Premium, Haiku 4.5, guideline removed & 32 & $+0.071$ & 0.119 & 0.333 & \textbf{no} \\
Premium, smaller models pooled & 128 & $+0.138$ & 0.114 & 0.318 & yes \\
GPT-6 Luna, structure-matched & 832 & $+0.220$ & 0.049 & 0.136 & yes \\
GPT-6 Luna, bare blocks & 832 & $+0.256$ & 0.049 & 0.138 & yes \\
GPT-6 Sol, structure-matched & 832 & $+0.239$ & 0.054 & 0.150 & yes \\
GPT-6 Sol, bare blocks & 832 & $+0.113$ & 0.044 & 0.123 & yes \\
Extra pull from removing guideline, Opus 5 & 320 & $+0.541$ & 0.071 & 0.199 & yes \\
Extra pull from removing guideline, Sonnet 5 & 64 & $+0.632$ & 0.152 & 0.426 & \textbf{no} \\
Extra pull from removing guideline, Haiku 4.5 & 64 & $+0.357$ & 0.063 & 0.176 & yes \\
\bottomrule
\end{tabular}
\end{table}


\section{Civil-damages results across model lineages}
\label{app:damages}

We applied the displaced-number design to sixteen fictional civil claims, crossing severity,
liability, mitigation, and documentation. D1 used Claude Opus~5. Its proximity measures were
selected after examining the first eight agents, so those panel estimates are exploratory; the
measures and direction were then registered before the two-agent confirmation sample (A12 and A01)
was examined. On those 32 matched pairs, in-band awards were 14/32 under the peer label and 26/32
under the tool label ($p=0.0042$, two-sided exact sign test), awards within two percent of the
central figure 10/32 and 24/32 ($p=0.00052$), and exact central matches 2/32 and 23/32
($p=9.54\times10^{-7}$). The peer-minus-tool difference in proportional distance was $+0.140$
(95\% $t$ $[0.032,0.248]$, MDE $0.142$). The registered rule calls this suggestive; see
\texttt{damages/AMENDMENT-D1-proximity.md}.

D2 was registered before collection and used the same ten agents, claims, primary arms, and
rendered prompts as the complete D1 panel; every prompt hash matched its D1 counterpart. GPT-6 Sol
at high reasoning returned 319 of 320 planned responses parseably (the A07 step-3 tool cell gave
none and was not retried), leaving 159 matched pairs. Table~\ref{tab:damages-cross} compares all
ten D1 agents with D2; its D1 column pools the exploratory panel with the confirmation sample and
is descriptive. Differences are peer-minus-tool with 95\% claim-clustered $t$ intervals; binary
$p$-values are exact sign tests on discordant pairs, reported peer-only/tool-only.

\begin{table}[H]
\centering
\caption{Civil-damages proximity outcomes by model run.}
\label{tab:damages-cross}
\scriptsize
\setlength{\tabcolsep}{3pt}
\renewcommand{\arraystretch}{1.0}
\begin{tabular}{>{\raggedright\arraybackslash}p{0.15\textwidth}>{\raggedright\arraybackslash}p{0.38\textwidth}>{\raggedright\arraybackslash}p{0.38\textwidth}}
\toprule
Outcome & Claude Opus 5, D1 (160 pairs) & GPT-6 Sol high, D2 (159 pairs) \\
\midrule
Inside band &
Peer 21/160; tool 88/160; $\Delta=-0.419$ [$-0.557,-0.281$]; $p=2.17\times10^{-18}$; discordant peer-only/tool-only $2/69$ &
Peer 90/159; tool 88/159; $\Delta=+0.013$ [$-0.127,+0.152$]; $p=0.871$; discordant $20/18$ \\
Distance from center &
Mean peer 0.392; tool 0.156; $\Delta=+0.236$ [$+0.145,+0.328$]; MDE 0.120 &
Mean peer 0.274; tool 0.182; $\Delta=+0.093$ [$+0.011,+0.174$]; MDE 0.107 \\
Within 2\% of center &
Peer 11/160; tool 68/160; $\Delta=-0.356$ [$-0.485,-0.227$]; $p=2.08\times10^{-16}$; discordant $1/58$ &
Peer 72/159; tool 73/159; $\Delta=-0.006$ [$-0.101,+0.089$]; $p=1.000$; discordant $17/18$ \\
Equals center &
Peer 2/160; tool 49/160; $\Delta=-0.294$ [$-0.390,-0.198$]; $p=1.42\times10^{-14}$; discordant $0/47$ &
Peer 34/159; tool 43/159; $\Delta=-0.057$ [$-0.152,+0.039$]; $p=0.200$; discordant $15/24$ \\
\bottomrule
\end{tabular}
\end{table}

The registered D2 primary outcome does not reproduce D1's large in-band difference: the contrast is
small, reverses direction, and is not statistically decisive. This is not evidence that the arms
are equivalent. For central distance, the D2 interval excludes zero, but its MDE exceeds the point
estimate, so the registered rule calls the result suggestive rather than decisive. Neither binary
secondary outcome rejects. The strong D1 proximity pattern therefore does not generalize unchanged
to this GPT-6 Sol sample.

For continuity, the pooled displacement slope is $+0.383$ on D1 (95\% $z$ interval
$[0.185,0.581]$, MDE $0.283$, within-agent permutation $p=0.0014$) and $+0.559$ on D2
(95\% $z$ interval $[0.344,0.774]$, MDE $0.307$, $p=0.0002$). The tool-minus-peer slope premium
is $+0.215$ on D1 (MDE $0.311$, $p=0.0712$) and $+0.232$ on D2 (MDE $0.292$, $p=0.0341$); both
estimates are below their MDEs. These secondary summaries do not replace the registered proximity
outcomes. D1 used Claude subagents, whereas D2 used Codex CLI, so the cross-run difference is
descriptive across both model and execution route. The results do not separate trust in a source
label from adoption of a supplied number.


\end{document}
